% =========================================================
% Proof: Preservation of Diameter by Isometries
% Source: notes/isometries/notes-isometries.tex
% Difficulty: Easy
% =========================================================

\subsection*{Preservation of Diameter}
\label{prf:isometry-preserves-diameter}

\begin{remark*}[Return]
$\leftarrow$ \hyperref[prop:isometry-preserves-diameter]{Back to Proposition in Notes}
\end{remark*}

\begin{proof}
\Claim If $f : X \to Y$ is an isometry and $A \subseteq X$, then
$\operatorname{diam}(f(A)) = \operatorname{diam}(A)$.

\Given An isometry $f$ and $A \subseteq X$.

\Goal To show $\sup\{d_Y(f(a),f(b)) : a,b \in A\} = \sup\{d_X(a,b) : a,b \in A\}$.

\Strategy The two distance-sets are equal because $d_Y(f(a),f(b)) = d_X(a,b)$
for every pair. Suprema of equal sets are equal.

% -------------------------------------------------------
% YOUR PROOF HERE
% -------------------------------------------------------

\end{proof}

\begin{remark*}[Proof shape]
The isometry condition converts every distance in $f(A)$ to the
corresponding distance in $A$; the distance sets are identical.
\end{remark*}

\begin{remark*}[Dependencies]
\begin{itemize}
  \item Definition~\ref{def:isometry}: $d_Y(f(x),f(y)) = d_X(x,y)$.
  \item Definition~\ref{def:diameter}: $\operatorname{diam}(A) = \sup\{d(x,y) : x,y \in A\}$.
\end{itemize}
\end{remark*}
